% Updated empirical section for the ZeBRA--genomics manuscript.
% Loaded at \begin{document} from empirical_figures.tex so it replaces the
% empirical-section macro defined in zebra_genomics_boundary_theorem_empirical.tex.

\renewcommand{\empiricalsection}{%
\origsection{Empirical correspondence in the Colorado ILD analysis}
\label{sec:empirical}

\subsection{Cohort, endpoint geometry, and analysis sequence}
The current analysis uses the Colorado genomic--clinical cohort assembled in \texttt{zebra\_comp}. The three literal stored phenotype fields are \texttt{FILD or FILA ADJUDICATED}, \texttt{NONFIBROTIC ILD ADJUDICATED}, and \texttt{NON FIBROTIC ILA ADJUDICATED}; standardized display labels are used below. The operational notebook target maps \texttt{Y/y} to 1 and \texttt{N/n} to 0, converts the field to numeric, assigns missing target values to 0, and excludes observations without an available 104-week ZeBRA score. The reported endpoint is therefore the notebook-defined operational prediction target; not every target-zero observation is an independently adjudicated negative. This distinction is retained explicitly when interpreting false-positive rates.

The processed genomic matrix contains 171 selected variant loci represented by categorical genotype states; the \textit{MUC5B} promoter variant rs35705950 is included explicitly. The empirical sequence mirrors the theorem. Notebook~32 asks whether genomics alone or a broad nonlinear ZeBRA+genomics model improves global held-out discrimination. Notebook~33 then asks the paired incremental-value question conditional on ZeBRA. Separate analyses test whether ZeBRA reconstructs \textit{MUC5B}, where residual \textit{MUC5B} information lies in ZeBRA score space, whether that local information can improve a boundary-adjacent operating decision, and how the channels behave at stringent FPRs.

\begin{table*}[!t]
\caption{Theory--evidence map. The theorem is mathematical; the Colorado results are correspondence tests and illustrations rather than proofs of A1--A2.}
\label{tab:theory-evidence-map}
\centering\footnotesize
\setlength{\tabcolsep}{6pt}
\begin{tabular}{p{0.23\textwidth}p{0.32\textwidth}p{0.36\textwidth}}
\toprule
Theoretical statement & Empirical evidence & Interpretation\\
\midrule
Distinct clinical and genomic evidence channels & ZeBRA predicts rs35705950 carrier status at chance; pooled tail enrichment disappears within FILD strata (Fig.~\ref{fig:muc5b-apparent-stratified}) & Weak global increment is not explained by implicit reconstruction of the major genotype.\\
Weak global genomic increment after a strong clinical score & Notebook-32 repeated-split comparison and notebook-33 paired conditional analysis (Figs.~\ref{fig:global-repeated-boxplots}--\ref{fig:incremental-auc}) & Under the evaluated genomic representations and learners, broad genomic augmentation supplies no reproducible global FILD/FILA gain once ZeBRA is known.\\
Boundary-localized residual genomic information & FPR-band and moving-window conditional-information analyses (Fig.~\ref{fig:local-information}; Tables~\ref{tab:local-info}--\ref{tab:local-windows}) & Residual \textit{MUC5B} information is concentrated in restricted clinical-score regimes rather than uniformly distributed.\\
Selective genomic information can change boundary-adjacent decisions & Leakage-free matched-operational-FPR rescue analysis (Fig.~\ref{fig:boundary-rescue}, Table~\ref{tab:boundary-selected}) & Targeted use of \textit{MUC5B} improves held-out sensitivity while preserving approximately the same operational FPR.\\
Clinical-tail dominance at stringent operating points & Direct held-out LR$+$ comparison at requested FPRs (Fig.~\ref{fig:low-fpr-geometry}) & ZeBRA has much larger LR$+$ than the evaluated genomic-only model at stringent operating points; this is correspondence to, not proof of, the bounded-LR theorem.\\
\bottomrule
\end{tabular}
\end{table*}

\subsection{Global discrimination: clinical history dominates}
The principal global comparison uses 30 repeated stratified outer splits with 40\% training and 60\% held-out evaluation. Within each split, ZeBRA is evaluated as the existing clinical-history score, a genomic-only LightGBM model is fitted to the genomic variables, and a combined LightGBM model is fitted to genomics plus ZeBRA. Hyperparameter selection occurs inside the training portion of each split.

For FILD/FILA, mean held-out AUC is 0.8142 for ZeBRA, 0.5662 for the genomic model, and 0.7988 for the nonlinear combined model (Fig.~\ref{fig:global-repeated-boxplots}; Table~\ref{tab:auc-repeated}). The combined model is below ZeBRA in all 30 paired splits, with mean combined-minus-ZeBRA AUC $-0.01535$. The two nonfibrotic outcomes show the same broad ordering but contain only 74 and 27 positive cases, respectively, and are treated as ancillary contrasts rather than strong phenotype replications.

\begin{table*}[!t]
\caption{Held-out AUC over 30 repeated outer splits from notebook~32. Values are mean $\pm$ SD.}
\label{tab:auc-repeated}
\centering\footnotesize
\begin{tabular}{lccc}
\toprule
Target & ZeBRA & Genomic & Combined\\
\midrule
FILD/FILA & $0.8142\pm0.0089$ & $0.5662\pm0.0245$ & $0.7988\pm0.0150$\\
Nonfibrotic ILA & $0.8281\pm0.0162$ & $0.4717\pm0.0443$ & $0.7112\pm0.0541$\\
Nonfibrotic ILD & $0.8202\pm0.0257$ & $0.5141\pm0.0719$ & $0.6145\pm0.0953$\\
\bottomrule
\end{tabular}
\end{table*}
\figGlobalRepeatedSplits

\subsection{Paired incremental-value analysis: essentially zero global FILD/FILA gain}
Because the nonlinear comparison alone cannot distinguish lack of information from a poor fusion architecture, notebook~33 compares a ZeBRA-only logistic model with a regularized elastic-net model containing ZeBRA plus the complete genomic panel on the same held-out subjects. For FILD/FILA, the paired mean AUC increment is only $+0.000435$, the median increment is exactly 0, and the split-wise 2.5th--97.5th percentiles are $[-0.00507,0.00408]$ (Fig.~\ref{fig:incremental-auc}). These are empirical split quantiles, not a classical confidence interval from independent samples.

The proper scoring rules point in the same direction. Adding the genomic panel decreases average precision by 0.0311 on average and worsens Brier score and log loss in all 30 FILD/FILA splits. Thus the supported claim is deliberately model-specific and predictive: under the evaluated genomic representations and augmentation models, broad genomics provides no reproducible global predictive gain once ZeBRA is included. It is not a claim of information-theoretic conditional independence.

\figIncrementalAUC
\begin{table*}[!t]
\caption{Paired incremental effect of adding the genomic panel to ZeBRA in notebook~33. Positive $\Delta$Brier and $\Delta$log-loss are worse.}
\label{tab:incremental-metrics}
\centering\scriptsize
\begin{tabular}{lrrrrrr}
\toprule
Target & Mean $\Delta$AUC & Median $\Delta$AUC & 2.5--97.5\% & Mean $\Delta$AP & Mean $\Delta$Brier & Mean $\Delta$log-loss\\
\midrule
FILD/FILA & +0.00043 & 0.00000 & $[-0.00507,0.00408]$ & -0.03111 & +0.000312 & +0.004680\\
Nonfibrotic ILA & -0.03737 & -0.03667 & $[-0.07384,-0.00705]$ & -0.01392 & +0.000040 & +0.000980\\
Nonfibrotic ILD & -0.10052 & -0.08615 & $[-0.22577,-0.00506]$ & -0.00847 & +0.000062 & +0.000973\\
\bottomrule
\end{tabular}
\end{table*}

\subsection{The null global increment is not explained by genotype reconstruction}
Among 12,766 subjects with called rs35705950 genotype, ZeBRA predicts T-carrier status at essentially chance level: pooled continuous-score AUC 0.5088, explicitly adjudicated FILD/FILA-negative AUC 0.4847, and FILD/FILA-positive AUC 0.4817 (Fig.~\ref{fig:muc5b-apparent-stratified}). Carrier prevalence is nevertheless 39.3\% in FILD/FILA-positive subjects versus 19.2\% among explicitly adjudicated negatives. A pooled top-1\% ZeBRA-tail carrier enrichment of 1.47-fold ($p\approx0.009$ by the repository permutation analysis) disappears after disease stratification. This is consistent with both channels being associated with disease status rather than with ZeBRA reconstructing genotype; conditioning on disease is not presented as a formal mediation analysis.

\figMUCfiveBApparentAndStratified

\subsection{Residual genomic information is localized in ZeBRA score space}
The absence of global incremental AUC does not imply zero conditional genomic information everywhere. We therefore examined both clinically interpretable control-FPR bands and overlapping moving windows along the ZeBRA percentile axis. These analyses test a qualitative implication of Theorem~\ref{thm:main}(ii): residual genomic information should be concentrated in restricted clinical-evidence regimes rather than distributed uniformly. They do not estimate the theorem's exact LR interval $[\lambda/M,\lambda/m]$, and the moving windows overlap and are not independent confirmatory tests.

The FPR-defined analysis makes the channel asymmetry explicit (Fig.~\ref{fig:local-information}A). In every displayed band, the LRT for adding \textit{MUC5B} after ZeBRA ($G\mid Z$) exceeds the converse LRT for adding ZeBRA after \textit{MUC5B} ($Z\mid G$). The largest differences occur in the 5--10\% band ($G\mid Z=16.55$ versus $Z\mid G=6.84$) and the 20--40\% band (13.69 versus 0.064). The nominal $G\mid Z$ p-values are $4.74\times10^{-5}$ and $2.15\times10^{-4}$, respectively. These p-values localize signal but are not multiplicity-adjusted confirmatory evidence.

The continuous moving-window analysis makes the regime transition particularly transparent (Fig.~\ref{fig:local-information}B; Table~\ref{tab:local-windows}). Across much of the intermediate-to-high ZeBRA axis, residual \textit{MUC5B} information substantially exceeds residual ZeBRA information. At the 70th-percentile window, $G\mid Z=10.93$ while $Z\mid G\approx0.002$; at the 90th percentile the corresponding values are 19.78 and 1.23; and at 92.5 they are 20.37 and 0.41. The ordering reverses only in the extreme clinical-risk tail: at the 95th-percentile window, $Z\mid G=58.76$ versus $G\mid Z=22.86$, and at 97.5 the values are 68.12 versus 15.61. The data therefore exhibit a local information crossover: genomics retains conditional predictive information through a broad boundary-adjacent region, whereas the clinical-history channel becomes dominant at the most extreme ZeBRA values. This is qualitative correspondence to the theorem's geometry, not an empirical proof of A1 or A2.

Panels C--D show the same pattern in directly interpretable quantities. Within ZeBRA windows, T-carriers retain higher observed FILD/FILA prevalence than non-carriers: 4.28\% versus 0.98\% at the 70th-percentile window, 8.27\% versus 2.05\% at 90, 18.45\% versus 7.75\% at 95, and 20.48\% versus 9.89\% at 97.5. The adjusted \textit{MUC5B} carrier OR remains several-fold across much of the risk axis, while the ZeBRA OR per local SD rises sharply in the extreme tail, reaching 2.95 at the 95th percentile and 4.07 at 97.5. These representations make the central empirical claim precise: the genomic signal is neither globally redundant nor uniformly useful; it is locally informative until sufficiently strong clinical evidence becomes dominant.

\figLocalInformation

\begin{table*}[!t]
\caption{Local FILD/FILA information in ZeBRA-defined control-FPR bands. Both directional conditional LRTs are shown.}
\label{tab:local-info}
\centering\scriptsize
\begin{tabular}{lrrrrrrrr}
\toprule
Band & $N$ & Cases & LRT $Z\mid G$ & LRT $G\mid Z$ & $G{-}Z$ partial & nominal $p(G\mid Z)$ & AUC ZeBRA & AUC \textit{MUC5B}\\
\midrule
2--5\% & 401 & 26 & 1.15 & 4.44 & 3.29 & 0.0352 & 0.577 & 0.595\\
5--10\% & 652 & 26 & 6.84 & 16.55 & 9.71 & $4.74\times10^{-5}$ & 0.367 & 0.677\\
10--20\% & 1,281 & 30 & 1.19 & 4.77 & 3.58 & 0.0289 & 0.558 & 0.586\\
20--40\% & 2,540 & 37 & 0.06 & 13.69 & 13.63 & $2.15\times10^{-4}$ & 0.516 & 0.635\\
\bottomrule
\end{tabular}
\end{table*}

\begin{table*}[!t]
\caption{Selected overlapping moving windows illustrating the conditional-information crossover.}
\label{tab:local-windows}
\centering\scriptsize
\begin{tabular}{rrrrrrr}
\toprule
ZeBRA percentile & LRT $Z\mid G$ & LRT $G\mid Z$ & OR ZeBRA/local SD & OR \textit{MUC5B} & Prev. carrier & Prev. non-carrier\\
\midrule
65.0 & 0.45 & 13.36 & 1.17 & 5.53 & 0.0466 & 0.00865\\
70.0 & 0.002 & 10.93 & 0.99 & 4.52 & 0.0428 & 0.00980\\
90.0 & 1.23 & 19.78 & 1.20 & 4.31 & 0.0827 & 0.0205\\
92.5 & 0.41 & 20.37 & 1.09 & 3.56 & 0.1067 & 0.0323\\
95.0 & 58.76 & 22.86 & 2.95 & 2.70 & 0.1845 & 0.0775\\
97.5 & 68.12 & 15.61 & 4.07 & 2.43 & 0.2048 & 0.0989\\
\bottomrule
\end{tabular}
\end{table*}

The unusually low within-band ZeBRA AUC of 0.367 in the 5--10\% control-FPR band is the direct source result, not a transcription error. It means that within that restricted slice, the raw ZeBRA score is not locally monotone for case ranking; it is another reason not to identify the uncalibrated raw score numerically with the theorem's clinical likelihood ratio.

\subsection{Boundary-targeted \textit{MUC5B} use improves held-out sensitivity at matched operational FPR}
The two-threshold rescue rule is evaluated leakage-free at the split level: upper and lower ZeBRA thresholds are derived from training negatives, the \textit{MUC5B} rescue is applied to untouched outer-test subjects, and a ZeBRA-only comparator is separately chosen on training data to match the rescue rule's training-set FPR. Across all six evaluated threshold pairs, the mean held-out sensitivity difference is positive (Fig.~\ref{fig:boundary-rescue}; Table~\ref{tab:boundary-selected}). For the 0.5\%/10\% pair, matched-ZeBRA sensitivity is 0.3587 and rescue sensitivity is 0.4033, a mean gain of 0.0446, while operational FPR changes from 0.02335 to 0.02310. Because six pairs were evaluated, the 4.46-point value is best treated as the maximum observed within this prespecified grid rather than as a single confirmatory effect estimate.

\figBoundaryRescue
\begin{table*}[!t]
\caption{All six evaluated matched-operational-FPR rescue rules.}
\label{tab:boundary-selected}
\centering\scriptsize
\begin{tabular}{lrrrrrrr}
\toprule
Upper/lower target & Matched ZeBRA sens. & Rescue sens. & $\Delta$ sens. & Matched FPR & Rescue FPR & Matched LR$+$ & Rescue LR$+$\\
\midrule
0.5\% / 2\% & 0.2919 & 0.3101 & +0.0181 & 0.00801 & 0.00802 & 38.16 & 39.67\\
0.5\% / 5\% & 0.3121 & 0.3487 & +0.0366 & 0.01404 & 0.01400 & 22.74 & 25.21\\
0.5\% / 10\% & 0.3587 & 0.4033 & +0.0446 & 0.02335 & 0.02310 & 15.58 & 17.54\\
1\% / 2\% & 0.3046 & 0.3142 & +0.0096 & 0.01227 & 0.01228 & 25.43 & 26.03\\
1\% / 5\% & 0.3303 & 0.3528 & +0.0225 & 0.01844 & 0.01826 & 18.19 & 19.52\\
1\% / 10\% & 0.3740 & 0.4075 & +0.0335 & 0.02753 & 0.02735 & 13.75 & 14.95\\
\bottomrule
\end{tabular}
\end{table*}

The operational target-zero group includes unadjudicated subjects. In the small explicitly adjudicated-negative subset, rescue FPR is higher than the original stringent baseline, but the current result set does not provide the corresponding explicitly-adjudicated-negative FPR for the fair matched-ZeBRA comparator. We therefore use the term \emph{matched operational FPR}; the adjudicated-negative matched comparison remains an additional analysis rather than a settled clinical-specificity claim.

\subsection{Stringent operating points directly support the clinical-tail correspondence}
The most direct empirical counterpart to Theorem~\ref{thm:main}(iii)--(iv) is the held-out LR$+$ comparison at stringent requested FPRs, not global AUC. At requested FPR 0.5\%, mean held-out LR$+$ is approximately 57.9 for ZeBRA, 4.28 for the genomic-only model, and 39.1 for the nonlinear combined model. At requested FPR 1\%, the corresponding means are approximately 29.8, 3.62, and 26.7. Figure~\ref{fig:low-fpr-geometry} shows the full requested-FPR series together with sensitivity. Thus the evaluated clinical-history score dominates the evaluated genomic-only model in the stringent tail. This does not prove a universal genomic bound $M_G$ or an unbounded ZeBRA LR tail, but it is the empirical quantity most closely aligned with the theorem's extreme-specificity statement.

\figLowFPRGeometry

\subsection{Calibration delimits the present theory--data mapping}
The field historically named \texttt{predicted\_risk} is not calibrated as an absolute FILD/FILA probability in this cohort: mean raw ZeBRA score is approximately 0.59 at an operational prevalence of approximately 0.020, with calibration slope approximately 0.18 and expected calibration error approximately 0.57. The theorem is therefore not applied numerically by treating raw ZeBRA as a posterior probability or likelihood ratio. In the present manuscript ZeBRA is an ordinal clinical-history score; the local and stringent-operating analyses provide qualitative correspondence. A stronger theory test should calibrate or directly estimate a score-level empirical likelihood ratio before comparing the observed transition with the theoretical $[\lambda/M,\lambda/m]$ band.

\subsection{Empirical synthesis and limitations}
The complete analysis supports a specific regime-dependent claim. First, ZeBRA is much stronger globally than the evaluated genomic panel. Second, broad genomic augmentation provides essentially no reproducible global FILD/FILA gain after ZeBRA is known. Third, that null global increment is not explained by ZeBRA reconstructing rs35705950. Fourth, both FPR-defined and continuous moving-window analyses show substantial residual \textit{MUC5B} information through intermediate/high ZeBRA regimes, with a direct crossover to clinical-history dominance only in the extreme ZeBRA tail. Fifth, selective \textit{MUC5B} use near an operating boundary improves held-out sensitivity at approximately matched operational FPR. Finally, direct stringent-FPR LR$+$ comparisons show that the evaluated clinical channel is far stronger than the genomic-only model in the high-specificity tail.

Several qualifications remain essential. A1 and A2 are mathematical regularity assumptions and are not established by the Colorado data. The selected genomic panel is not the strongest possible genomic comparator; an explicit \textit{MUC5B}-only baseline and an externally specified contemporary IPF/fibrotic-ILD PRS would strengthen the comparison. The moving windows overlap, and the local tests and rescue grid are exploratory within the same Colorado resource. The two nonfibrotic outcomes are sparse ancillary contrasts. Repeated train/test splits share subjects and their quantiles are not independent-sample confidence intervals. The operational endpoint maps missing adjudication to zero. Independent or prospectively prespecified replication is therefore required before treating the local rescue rule as a clinical decision rule.
}
